\chapter{Inversion}
\label{chap:inversion}

% Closed question: How can information move between histories while preserving
% admissibility?

\begin{chapteressay}

The previous chapter ended with a difficult conclusion: the selection
rule that picks one admissible continuation from many is provably
uncomputable. The selection rule requires exact Kolmogorov
complexity; Kolmogorov complexity requires the Busy Beaver function;
the Busy Beaver is not computable. The chapter ended with the
observation that the universe is doing something no Turing machine
can simulate, while the compiler is restricted to proxies.

This chapter continues the algorithmic vocabulary. The previous
chapter's noncomputability did not stop the algorithm from doing
useful work. Many algorithms still operate admissibly, even though
the selection rule itself is uncomputable. This chapter examines
one specific class of admissible algorithms: those that move
information between histories through inversion.

Inversion is the algorithmic operation of taking a forward
computation and computing its adjoint. The forward computation
goes from inputs to outputs; the adjoint goes from outputs back
to inputs, exposing the sensitivity of the output to each input.
The two computations are related by an algebraic identity: the
adjoint of the linearization. The identity is the chain rule,
applied to the entire algorithm.

The classical example is matrix algebra. The system $Ax = b$ has
solution $x = A^{-1} b$. Direct computation of $A^{-1}$ would be
expensive; the standard algorithm is LU decomposition: $A = LU$
with $L$ lower triangular and $U$ upper triangular. The forward
substitution $Ly = b$ produces an intermediate; the backward
substitution $Ux = y$ produces the solution. The two substitutions
together solve the system in $O(n^3)$ for the decomposition and
$O(n^2)$ for the two substitutions.

The forward-backward pattern repeats throughout the chapter. In
neural networks, the forward pass computes the loss; the backward
pass (backpropagation) computes the gradient. In Gram-Schmidt,
the forward direction projects onto previous basis vectors; the
backward direction subtracts those projections, leaving the
orthogonal residue. In quantum mechanical alpha decay, the forward
direction is the wave function's evolution; the backward
extraction is the probability of tunneling through the barrier.
In Lean's typeclass machinery, the forward direction is
deduction; the \texttt{WITNESSED} certificate is the backward
verification.

The second concrete example for the chapter is an expense audit.
An employee submits a report with sixty-two line items totaling
several thousand dollars. The accounting clerk audits backward:
from the total, back to the line items, back to the supporting
receipts, back to the merchants where the money was spent. Each
line item is a transport: a transport of an expense from a
merchant to the report. The audit is the inversion: from the
report back to the merchants. Each line item requires four
witnesses: the source (the merchant), the process (the credit
card transaction), the truth (the supporting receipt), and the
admissibility (the policy allows this category at this amount).
Missing witnesses produce metavariables, which the audit
sends back to the employee for clarification.

The expense audit makes the chapter's structure concrete in a
non-computational setting. The same algorithmic pattern
(forward computation, backward inversion, structured witnessing)
appears in linear algebra (LU), machine learning (backprop),
classical algorithms (Gram-Schmidt), quantum mechanics (alpha
decay), and accounting (expense audits). The pattern is
algorithmic, not domain-specific. The chapter's claim is that
this pattern is the algorithmic form forced by the chapter's
question.

The chapter's closed question is:

\begin{center}
\emph{How can information move between histories while preserving
admissibility?}
\end{center}

The answer is: through inversion, with the
\texttt{WITNESSED} structure as the formal certificate. The
forward computation transports information; the inversion
provides the backward verification; the four-input witness
(source, process, truth, admissibility) certifies that the
transport was admissible.

The chapter has a particular relationship to Chapter 6. Chapter
6's argument was that the selection rule is uncomputable.
Chapter 7's argument is that, despite the uncomputability of the
selection, many admissible transports are computable. The
compiler can certify them. The proofs of admissibility are
finite. The \texttt{WITNESSED} certificates can be produced.

The book's overall argument is therefore not that all is lost. It
is that most operations are admissible and certifiable; only the
selection rule itself is the uncomputable boundary. The compiler
handles the bulk of the work. The exceptional case is when the
algorithm must select between near-tied admissible continuations,
and that selection requires the exact $K$ that no algorithm can
compute. For the bulk of practical inversion problems, the
compiler is enough. The chapter's algorithmic discipline is the
working machinery of the admissible.

Subsequent chapters build on this. Chapter 8 addresses the
calibration that makes multiple observers' records commensurable.
Chapter 9 addresses the commutation failure when two transports
do not produce the same result. Chapter 10 addresses the
invariants that survive across admissible recompilations. Chapter
11 addresses the append-only structure of the final ledger and
the type-check that closes the book.

\end{chapteressay}

\section{LU Decomposition}
\label{se:lu-decomposition}

A linear system $Ax = b$ asks: given a matrix $A$ and a vector $b$,
find the vector $x$ such that $Ax = b$. For a $1000 \times 1000$ matrix,
the naive solution (Cramer's rule, computing determinants) is
$O(n!)$ --- impossibly slow. The standard solution is LU
decomposition.

The matrix $A$ is factored as $A = LU$, where $L$ is lower triangular
(zero above the diagonal) and $U$ is upper triangular (zero below the
diagonal). The factorization is computed by Gaussian elimination,
which takes $O(n^3)$ operations: roughly $10^9$ multiplications for a
$1000 \times 1000$ matrix, which a modern CPU performs in a fraction
of a second.

Once the factorization is computed, the original system $Ax = b$
becomes $L(Ux) = b$, which factors into two triangular systems. Let
$y = Ux$. Then $Ly = b$ is solved by forward substitution (start with
$y_1$, work down). $Ux = y$ is solved by backward substitution (start
with $x_n$, work up). Each substitution is $O(n^2)$. The total cost
is $O(n^3 + 2n^2)$, dominated by the initial decomposition.

The two substitutions are forward and backward. Forward goes from the
input ($b$) to the intermediate ($y$). Backward goes from the
intermediate ($y$) to the output ($x$). The two sweeps are the
algorithm's traversal in both directions: forward to solve, backward
to invert.

The forward sweep is the canonical forward computation. Given the
matrix and the right-hand side, the algorithm produces the
intermediate. The backward sweep is the canonical adjoint
computation. Given the intermediate and the upper triangular factor,
the algorithm produces the answer. The backward sweep inverts the
forward sweep's action.

This is the chapter's question in linear algebra form. Information
flows forward (from $b$ to $y$ via $L$) and then backward (from $y$
to $x$ via $U$). The two flows together solve the system. Neither
flow alone produces the answer; both are needed. The algorithm
preserves admissibility at each step: every operation is a single
row reduction, every reduction is reversible, every reduction
preserves the system's solution set.

The Lean analog is the typeclass cascade in \texttt{Episode8.lean}.
The instance cascade \texttt{STEP\_1} through
\texttt{ACOLYTE\_PROPAGANDA} formalizes the algorithmic transport
of structure through a stack of typeclass instances. Each step in
the cascade is one inference rule application: forward propagation
of structure. The reverse application (the adjoint) propagates
constraints backward: from a desired output type, the algorithm
infers the required input type.

LU decomposition is the simplest example of forward and adjoint
operations in the same algorithm. The forward operation is
multiplication by the matrix: $Ax = b$. The adjoint operation is
multiplication by the transposed matrix: $A^T y = c$. The two
operations are related: if $A x = b$, then $\langle A x, y
\rangle = \langle b, y \rangle$, which by the adjoint property
equals $\langle x, A^T y \rangle$. The adjoint property is the
algebraic structure that allows the backward sweep to invert the
forward sweep.

This is the foundation of the chapter's notion of inversion.
Information moves through an algorithm via a sequence of admissible
transformations. The inversion of the algorithm moves information
backward through the same sequence. The forward and backward
movements are related by an algebraic identity (the adjoint
property). The identity is what guarantees that the inversion is
admissible: applying the forward and then the backward operation
returns to the original input.

\section{Backpropagation}
\label{se:backpropagation}

A five-layer neural network is trained to classify handwritten digits.
The input is a $28 \times 28$ image. The output is a probability
distribution over ten digit classes. The middle three layers have
varying widths --- say 512, 256, 128 neurons. The total number of
weights is roughly $10^6$. The training data is the MNIST dataset:
60,000 training images and 10,000 test images.

Training proceeds by gradient descent on the cross-entropy loss. For
each training image, the network produces a predicted distribution.
The loss is the cross-entropy between the predicted distribution and
the true label (a one-hot vector). The training adjusts the weights
to decrease the loss.

The gradient of the loss with respect to each weight is what tells
the algorithm which way to move each weight. Computing this gradient
naively, by perturbing each weight individually and observing the
change in loss, would require $10^6$ forward passes per training step,
which would be hopelessly slow.

The trick is backpropagation. Compute one forward pass: input the
image, propagate through each layer, output the predicted
distribution, compute the loss. This is the forward direction. Then
compute a backward pass: starting from the loss, propagate the
gradient backward through the layers, applying the chain rule. The
backward pass takes about the same time as the forward pass. The
gradients of all $10^6$ weights are computed in one backward pass.

The backward pass is the inversion. The forward pass produces the
loss from the inputs. The backward pass produces the gradients of
the loss with respect to each weight. The two passes together fit
the chapter's structure: forward computes the result, backward
computes the sensitivity of the result to changes in the input. The
sensitivity is information about the algorithm's structure that
propagates backward from the result to the parameters.

The chain rule is the algebraic identity that makes
backpropagation work. For a composition of functions $f \circ g$,
the gradient is $\nabla(f \circ g)(x) = J_g(x)^T \cdot \nabla
f(g(x))$, where $J_g$ is the Jacobian of $g$. The transposed
Jacobian is the adjoint of the linear approximation of $g$. The
gradient at $x$ is the gradient at $g(x)$ pulled back through
$g$'s adjoint.

This is the same adjoint structure as LU decomposition. The
forward computation is the composition; the backward computation
is the adjoint of the forward. The two are related by the chain
rule. Each layer's contribution to the gradient is computed by
applying the adjoint of that layer.

The chapter's question is how information moves backward
admissibly. In backpropagation, the gradient information flows
backward from the loss to the parameters. The flow is admissible
because each step is the adjoint of a forward step. The adjoint
preserves the linear algebraic structure of the forward
computation. The gradient at the start (the input layer) is
related to the gradient at the end (the loss) by the
composition of the adjoints of all the layers.

The Lean analog is the typeclass cascade
\texttt{Knowledge} $\to$ \texttt{ScientificProcess} $\to$
\texttt{TRUTH} in \texttt{Episode9.lean}. Each transition in the
cascade represents the transport of a fact backward from a
conclusion to its supporting evidence. The conclusion is the
forward output; the evidence is the backward result. The
\texttt{Knowledge.le} comparison routes through the
\texttt{ScientificProcess} backward, retrieving the chain of
inferences that produced the conclusion.

Backpropagation is the chapter's most spectacular example of
inversion at scale. A neural network with $10^9$ parameters can
be trained efficiently because the gradient of all parameters
is computable in one backward pass per training example. The
cost is $O(P)$ for $P$ parameters per pass, asymptotically
optimal. Without backpropagation, the cost would be $O(P^2)$,
which for $P = 10^9$ would be infeasible.

The success of deep learning depends on backpropagation. The
algorithm is the workhorse of modern AI. The chapter's claim
is that backpropagation is not an optional optimization; it is
the algorithmic form of admissible inversion. Any system that
inverts a multi-layer computation must use the adjoint
structure that backpropagation embodies. The structure is
forced by the chapter's question: how does information move
backward without inventing structure not in the forward
computation?

\section{Gram-Schmidt Orthogonalization}
\label{se:gram-schmidt-orthogonalization}

Given three linearly independent vectors in three-dimensional space ---
say $v_1 = (1, 0, 0)$, $v_2 = (1, 1, 0)$, $v_3 = (1, 1, 1)$ --- the
Gram-Schmidt procedure produces three orthonormal vectors that span
the same space. The procedure is iterative.

Step 1: $u_1 = v_1 / \|v_1\| = (1, 0, 0)$. The first vector is the
first input, normalized.

Step 2: Subtract the projection of $v_2$ onto $u_1$. The projection is
$\langle v_2, u_1 \rangle u_1 = 1 \cdot (1, 0, 0) = (1, 0, 0)$. The
remainder is $v_2 - (1, 0, 0) = (0, 1, 0)$. Normalize: $u_2 = (0, 1,
0)$.

Step 3: Subtract the projections of $v_3$ onto $u_1$ and $u_2$. The
projection onto $u_1$ is $1 \cdot (1, 0, 0) = (1, 0, 0)$; onto $u_2$
is $1 \cdot (0, 1, 0) = (0, 1, 0)$. The remainder is $v_3 - (1, 0, 0)
- (0, 1, 0) = (0, 0, 1)$. Normalize: $u_3 = (0, 0, 1)$.

The three output vectors $u_1, u_2, u_3$ are orthogonal (pairwise
inner products are zero) and normalized (each has unit norm). They
span the same space as the input vectors $v_1, v_2, v_3$. The
procedure has produced an orthonormal basis for the input span.

Each step is an inversion. The projection of $v_2$ onto $u_1$ is
the forward direction: given $v_2$ and $u_1$, compute the
component of $v_2$ in the $u_1$ direction. The subtraction is the
inversion: remove that component, leaving the orthogonal residue.
The residue is the part of $v_2$ that is admissible as a new basis
vector.

The chapter's question manifests here as: how does the next basis
vector get oriented while preserving its admissibility as part of
the orthonormal basis? The answer is: by inverting the projections
onto the previous basis vectors. Each previous component is removed
through an inversion. The residue is the new admissible direction.

The procedure has a numerical issue: if the inputs are nearly
linearly dependent (almost in the same direction), the residue
becomes small and floating-point error becomes significant. The
modified Gram-Schmidt procedure addresses this by subtracting one
projection at a time (rather than all at once), which improves
numerical stability. The chapter's algorithmic content is the
same; the implementation differs.

Gram-Schmidt is the algorithmic form of building an admissible
basis from a sequence of inputs. The construction is incremental:
each new vector is added after the previous ones have been
processed. The processing is the inversion of the projections.
The output basis is admissible at each stage: after the first
vector, the basis spans a one-dimensional subspace; after the
second, two dimensions; and so on.

The Lean analog appears in the typeclass machinery for
\texttt{Knowledge.le} in \texttt{Episode9.lean}. The order
relation on knowledge states is: one knowledge state refines
another when it contains all the conclusions of the other plus
possibly more. The Gram-Schmidt procedure is the algorithmic
form of building a knowledge state incrementally: each new
fact is added after removing its component in the previously
established facts. The remaining residue is the new admissible
contribution.

\section{Alpha-Decay Transport}
\label{se:alpha-decay-transport}

Uranium-238 has a half-life of $4.47 \times 10^9$ years. A single
U-238 nucleus, isolated in a vacuum, has a fifty percent probability
of decaying within this time. The decay produces an alpha particle (a
helium-4 nucleus, two protons and two neutrons) emitted from the U-238
nucleus, leaving behind thorium-234.

The mechanism is quantum mechanical tunneling. Inside the nucleus,
the alpha particle is bound by the strong nuclear force; outside the
nucleus, it is repelled by the electrostatic force of the remaining
protons. Between these two regimes is the Coulomb barrier: a
potential energy maximum that classically the alpha particle cannot
cross because its kinetic energy is less than the barrier height.

Quantum mechanically, the alpha particle's wave function has a small
but nonzero amplitude on the other side of the barrier. The
probability of tunneling per unit time is exponentially small in the
barrier's integrated height. For U-238, the tunneling rate is small
enough to give a 4.47-billion-year half-life. For other isotopes,
the rate varies: polonium-214 has a half-life of 164 microseconds;
bismuth-209 has a half-life of $2 \times 10^{19}$ years.

The Gamow factor is the exponential rate of tunneling. It is
computable from the alpha particle's energy and the Coulomb
barrier's shape:
\[
P_{\text{tunnel}} \propto \exp\left(-\frac{2}{\hbar}
\int_{r_1}^{r_2} \sqrt{2m(V(r) - E)} \, dr\right),
\]
where $r_1$ and $r_2$ are the classical turning points, $V(r)$ is
the potential energy, $E$ is the alpha particle's energy, $m$ is its
mass, and $\hbar$ is Planck's constant.

This is the quantum-mechanical analog of the chapter's question.
Information (the alpha particle) moves from inside the nucleus
(where it is admissible as a bound state) to outside (where it is
admissible as a free particle), through a region where neither
state is admissible classically. The transport is admissible because
the wave function's continuity carries the alpha particle through
the barrier.

The transport is not the chapter's adjoint structure --- it is not
a backward computation. It is a different kind of admissible
information movement: through a region the classical algorithm
cannot traverse. The quantum mechanics provides the algorithmic
form: the wave function's amplitude varies continuously, including
through the barrier. The probability of a transport event is the
square of the amplitude at the other side. The transport rate is
the rate at which this probability accumulates.

The chapter's claim is that this is the algorithmic form of
transport through a classically forbidden region. The forward
computation (the wave function evolution) is admissible at every
point, including inside the barrier. The probability of a
transport event (the integrated amplitude at the far side) is the
admissible record of how much information has crossed the
barrier.

Alpha-decay transport is the algorithmic version of a
\texttt{ReligiousProcess} from \texttt{Episode9.lean}. A
religious process is one that transports belief through a
boundary that admissible inference cannot cross directly. The
boundary is the gap between empirically verifiable facts and
unverifiable beliefs; the transport is the propagation of belief
through the boundary. The chapter's algorithmic claim is that
this transport is structurally similar to alpha decay: both
involve information moving through a region where direct
admissibility is suspended.

The \texttt{Gospel} type in \texttt{Episode9.lean} is the
typeclass for the content of a \texttt{ReligiousProcess}. A
\texttt{Gospel} is a structured statement that the
\texttt{ReligiousProcess} transports. The structure includes the
content (the claim being transported), the warrant (the basis
for the transport), and the consequence (what the receiver is
expected to do with the transported content). The typeclass
makes the structure explicit.

The Lean compiler does not, of course, simulate alpha decay. The
typeclass is satirical: it formalizes the algorithmic structure
of belief transport through unobservable mechanisms. The point of
the satire is that algorithms that handle this kind of transport
must declare their warrant explicitly. The transport is not
admissible by default; it requires the warrant to be carried
through.

The chapter's transport themes thus span multiple registers:
linear algebra (LU decomposition), machine learning
(backpropagation), classical algorithms (Gram-Schmidt), quantum
mechanics (alpha decay), and Lean typeclass machinery
(\texttt{Knowledge.le}, \texttt{ReligiousProcess}). In every
case, the transport is admissible only when the algorithm
provides a structural justification: the adjoint property, the
chain rule, the projection-removal, the wave function continuity,
or the typeclass warrant. The justification is part of the
transport, not separate from it.

\section{WITNESSED in Lean}
\label{se:witnessed-in-lean}

\texttt{Episode10.lean} defines the \texttt{WITNESSED} typeclass that
formalizes the transport of a truth through an admissible process to
the current context. The class has the shape:

\begin{leanexcerpt}{WITNESSED}
class WITNESSED
    (Value : Type)
    (Carrier : CarrierProcess Value)
    [DISTINGUISHABLE Value Carrier] ... [TRUTH Value Carrier]
  where
  baptism : ReligiousProcess Value Carrier
  witness : Gospel
  risen?  : Gospel $\to$ Gospel $\to$ Prop := fun a b =>
    Gospel.le a b $\to$ Gospel.le b a
\end{leanexcerpt}

A \texttt{WITNESSED} instance carries the full project cascade
through \texttt{TRUTH}: every prior class is a typeclass
requirement on the bracketed list. The structure provides a
\texttt{baptism} (a \texttt{ReligiousProcess} for the underlying
witness), a current \texttt{witness} (a \texttt{Gospel}, the
project's name for the chapter's transported truth), and a
binary decision rule \texttt{risen?} that certifies the
transported gospel survives the round-trip. The compiler enforces
every link in the cascade.

The structure is the formal version of the chapter's transport
discipline. To transport a truth from one context to another, the
algorithm must provide: the source (what was true elsewhere), the
warrant (the process that carried it), the verification (that the
source was indeed true), and the admissibility (that the
destination is compatible). The four inputs together constitute the
witness for $p$.

Without the four inputs, $p$ cannot be witnessed in the current
context. The compiler enforces this: a use of $p$ that requires a
\texttt{WITNESSED} instance must provide all four. If any one is
missing, the compiler refuses to admit the use. The discipline is
strict: transport without witnessing is inadmissible.

The chapter's algorithmic claim is that \texttt{WITNESSED} is the
typeclass formalization of transport's admissibility. Every
admissible transport produces a \texttt{WITNESSED} certificate.
Every \texttt{WITNESSED} certificate requires the four inputs:
source, process, truth, admissibility. The four inputs make the
transport admissible; their absence makes it inadmissible.

The four inputs match the chapter's algorithmic examples. The
source is the input data (the right-hand side $b$ in LU, the
loss in backpropagation, the initial vector in Gram-Schmidt, the
bound state in alpha decay). The process is the forward
computation (the matrix factorization, the network's forward
pass, the orthogonalization, the wave function evolution). The
truth is the verification that the source has the claimed
property (the matrix's invertibility, the loss function's
differentiability, the input vectors' independence, the bound
state's energy). The admissibility is the relation between
source and destination (the solution belonging to the same vector
space, the gradient corresponding to the same parameter space,
the orthogonal basis spanning the same subspace, the free
particle representing the same nucleus).

The Lean machinery captures all four. The compiler's enforcement
ensures that every transport in user code provides all four. The
compiler is implementing the chapter's algorithmic discipline at
elaboration time. A user who attempts to transport a truth
without a \texttt{WITNESSED} certificate gets a compilation
error. A user who provides the certificate gets the transport
admitted.

The chapter's bridge to subsequent chapters: transport is now
algorithmically certified. The next question is how multiple
transports compose --- how distinct observers, each carrying their
own \texttt{WITNESSED} certificates, can compare their records.
Chapter 8 takes up the question of calibration: how the
reference structure is built such that two observers' counts can
be compared as one physical record.

\begin{bridge}

The chapter's question was how information can move between histories
while preserving admissibility.

The answer is inversion. The forward computation produces a result;
the inversion produces the sensitivity of the result to the inputs.
LU decomposition factors a matrix into forward (lower) and backward
(upper) triangular parts. Backpropagation computes gradients by
applying the chain rule backward through a neural network's layers.
Gram-Schmidt builds an orthonormal basis by inverting projections.
Alpha-decay transport carries information through a classically
forbidden region via wave function continuity. The
\texttt{WITNESSED} typeclass in Lean certifies that the transport
is admissible by requiring source, process, truth, and
admissibility.

What remains is the calibration that makes multiple transports
commensurable. Two observers, each carrying their own
\texttt{WITNESSED} certificates, must share a reference structure to
compare their records. Chapter 8 develops the ruler algorithm: the
calibration as hidden reference scale, the metric as the data
structure for frame-comparison.

\end{bridge}

\begin{coda}{Auditing an Expense Report Backward}

The expense report arrives at the corporate accounting department.
The employee submitting it has been traveling for a week. The report
is fourteen pages, with sixty-two line items, totaling four thousand
two hundred and eighteen dollars and sixteen cents. The cover page
lists the trip dates, the destinations, the project code, the
manager's approval signature, and the employee's certification that
the expenses are accurate.

The accounting clerk who receives the report has been at the
department for six years. She processes about two hundred expense
reports per week. Her job is to verify each report against the
company's expense policy and against the supporting receipts. If
everything checks out, she approves the report and routes it to
payroll for reimbursement. If something is wrong, she sends the
report back to the employee with notes about the discrepancies.

Today's report has a flag. The system has noted that the total
amount is in the ninety-fifth percentile for this employee's
historical pattern. The flag is a routine check; it triggers a
closer audit but does not imply wrongdoing.

The audit proceeds backward.

She starts with the total amount on the cover page: four thousand
two hundred and eighteen dollars. The total is a sum. She verifies
the sum by re-adding the line items. The line items total four
thousand two hundred and eighteen dollars and sixteen cents,
matching the cover page. The arithmetic is correct.

She moves backward to the line items. There are sixty-two of them,
each with a date, a description, a category code, and an amount.
She works through them in reverse order, starting with the most
recent (the employee's flight home, three days ago). The flight is
seven hundred forty dollars, charged to the project code, with a
boarding pass attached. The boarding pass shows the employee's
name, the flight number, the date, and the destination. The
boarding pass matches the report. She approves the line item.

She continues backward. The hotel for the four nights before the
return flight is one thousand two hundred dollars: three hundred
per night, with the hotel's receipt attached. The receipt shows
the hotel name, the dates, the room rate, and the taxes. The
receipt matches. The dates match the trip's authorized period.
She approves the line item.

She reaches a line item for an expensive dinner: two hundred
sixty-five dollars at a restaurant. The receipt is attached. The
date is the second day of the trip. The category code is
``client entertainment.'' She checks the company's policy:
client entertainment is allowed up to fifty dollars per person
per meal. The receipt does not specify the number of attendees.
She makes a note: the receipt requires clarification.

The audit's backward direction is the chapter's adjoint
operation. The forward direction is the employee's spending:
each purchase produces a receipt, the receipts accumulate into
the report, the report is submitted. The backward direction is
the clerk's verification: each line item is examined, the
supporting receipt is checked, the policy is verified, the
arithmetic is confirmed. The audit is the inversion of the
spending: from the total back to the individual transactions.

The chapter's \texttt{WITNESSED} structure appears at each step.
Each line item is a transport: a transport of an expense from
the merchant (the source) through the credit card transaction
(the process) to the report (the destination). The supporting
receipt is the \texttt{TRUTH} certificate; the matching of the
receipt to the line item is the admissibility check. The clerk's
approval is the witness: the certification that the transport
was admissible.

If a line item lacks a supporting receipt, the witness cannot be
produced. The clerk cannot certify the line item as transported
admissibly. The line item is denied; the employee must produce
the receipt or accept that the expense is not reimbursable. The
denial is the missing witness, recorded in the audit trail.

The dinner receipt that requires clarification is a metavariable:
the line item is constrained but not yet resolved. The clerk
notes the metavariable and continues the audit. At the end, she
will compile the list of metavariables and send it back to the
employee. The employee will provide the clarification, the
metavariables will resolve, and the audit can complete.

The clerk reaches the early line items, the ones from the start
of the trip. A taxi from the airport to the hotel on day one.
Twenty-eight dollars. The receipt shows the date, the
destination, the fare. The receipt matches. She approves it.

She has worked through all sixty-two line items. Five have
clarification requests. The remaining fifty-seven are approved.
The total approved is three thousand nine hundred fifty-three
dollars. The remaining two hundred sixty-five dollars (the
dinner) is held pending clarification.

She writes a memo to the employee. The memo lists the five
clarification requests, each with the line item number, the date,
the amount, and the specific question (number of attendees for
client entertainment; conversion rate for foreign currency
charges; receipt for a taxi without a receipt; project code for
an item that lacked one; date verification for an item with an
ambiguous date stamp).

The memo is the algorithmic inversion's record: it propagates
the audit's findings backward to the source (the employee) for
resolution. The employee will receive the memo, gather the
clarifying information, and re-submit. The clerk will then
complete the audit.

This is the chapter's structure in administrative form. Each
expense is a transport: from a merchant, through a payment, to a
report. The audit is the inversion: from the report, back through
the payments, to the merchants. The supporting receipts are the
\texttt{TRUTH} certificates that warrant each transport. The
matches between receipts and line items are the
\texttt{Knowledge.le} comparisons that establish admissibility.
The clerk's approvals are the \texttt{WITNESSED} certificates.

At the end of the day, the clerk closes the audit. The report is
partially approved (the fifty-seven items totaling three thousand
nine hundred fifty-three dollars are routed to payroll). The
partial approval is the admissible state: the report is not
finalized, but the approved portion is releasable. The five
pending items remain in the system, awaiting clarification.

The clerk's audit trail records every action: which items were
approved, which were flagged, which receipts were checked, which
questions were asked. The trail is the chapter's algorithmic
record of the inversion. Every transport (line item) has been
either approved (with the four witnesses) or flagged (with the
metavariable identifying what is missing). The audit is
complete in the algorithmic sense: every line item has a
status.

Tomorrow the clerk will process more expense reports. The
process will repeat. Each report is a forward computation; each
audit is the backward inversion. The system handles thousands of
such pairings per month. The algorithmic discipline is what
makes the system trustworthy: every dollar that leaves the
company has a paper trail, and every paper trail is auditable
in the backward direction.

\end{coda}
